%قبل از شروع مسئله ترجیحاً چیزی ننویسید. اما اگر خواستید بنویسید، در قالب لیستی (فقرات و شمارش) بنویسید.
%برای اضافه کردن نمره، بعدی از شروع‌مسئله، [x] را اضافه کنید و به جای x  نمره مورد نظر را قرار دهید. برای زیرمسئله‌ها هم به همین صورت است.
%\شروع‌مسئله[10]
\شروع‌مسئله
% صورت اصلی مسئله را در این‌جا بنویسید.
در یک آرایه با اندازه $n$، به عضوی که بیشتر از $ \frac{n}{k} $ بار تکرار شده‌است، عضو اکثریت می‌گوییم. راهکاری ارائه دهید که با پیچیگی حافظه 
$O(k)$
عضو اکثریت را در صورت وجود پیدا کند.\\
	\href{https://www.geeksforgeeks.org/majority-element/}{
[راهنمایی : راه حل برای
 $k = 2$]}

%اگر مسئله چند قسمت داشت به این صورت عمل کنید:
%\شروع{زیرمسائل}
%\زیرمسئله[0] % مسئلۀ ۱آ
%\زیرمسئله % مسئلۀ ۱ب
%\زیرمسئله % مسئلۀ ۱پ
%\زیرمسئله % مسئلۀ ۱ت
%\پایان{زیرمسائل}
\solution{
برای حل این مسئله، 
$k-1$
تا زوج مرتب تعریف می‌کنیم که مولفه‌ی اول هر زوج مرتب، نشان‌دهنده‌ی یک عدد است و مولفه‌ی دوم، نشان‌دهنده‌ی تکرر آن عدد است. فرض می‌کنیم که آرایه به شکل استریم ورودی گرفته می‌شود.

در هر گام، به شکل زیر عمل می‌کنیم. در ابتدا تمام زوج مرتب‌ها خالی است.

\begin{itemize}
\item
یک خانه از آرایه را ورودی می‌گیریم.
\item
اگر یکی از زوج مرتب‌ها بود که مولفه‌ی اولش با مقدار خوانده شده یکی بود، تکرر آن را یکی زیاد می‌کنیم.
\item
اگر یکی از زوج مرتب‌ها بود که خالی بود، مقدار آن را برابر با مقدار خوانده شده قرار می‌دهیم و تکرر آن  را یک می‌کنیم.
\item
در غیر این‌صورت، تکرر تمام زوج‌مرتب‌ها را یکی کم می‌کنیم. اگر تکرر یک زوج مرتب ۰ شد، فرض می‌کنیم که آن زوج مرتب خالی شده است.
\end{itemize}
حال ادعا می‌کنیم که اگر یک عضو اکثریت وجود داشته باشد، حتما مولفه‌ی اول یکی از زوج مرتب‌ها خواهد بود. 

اثبات: فرض کنید که یک مقدار اکثریت وجود داشته باشد. هر بار که یک عدد در ورودی می‌خوانیم، یکی از دو حالت زیر اتفاق می‌افتد.
\begin{itemize}
\item
تکرر یک زوج مرتب را افزایش می‌دهد.
\item
از تکرر تمام زوج‌مرتب ها یکی کم کند.
\end{itemize}
می‌دانیم که در مجموع حداکثر n بار تکرر یک زوج مرتب افزایش می‌یابد. چون اتفاق دوم، هر بار k تا از مجموع تکرر اعداد کم می‌کند، حداکثر 
$\frac{n}{k}$
بار اتفاق می‌افتد. حال فرض کنید که برای عضو اکثریت، a بار اتفاق اول افتاده باشد و b بار هم اتفاق دوم. می‌دانیم :
$$a + b > \frac{n}{k} \rightarrow a > \frac{n}{k} - b$$
پس تعداد دفعاتی که تکرر عدد اکثریت افزایش می‌یابد، بیشتر از تعداد دفعاتی است که تکررش کم می‌شود. بنابراین حتما این عضو در آخر کار، مولفه‌ی اول یک زوج مرتب است.

حال برای این که مطمئن شویم این عضو، یک عضو اکثریت است، کافی است یک بار دیگر تعداد تکرار های این عضو مشخص را در آٰرایه بشماریم(یعنی نیازمند این هستیم که آرایه دوبار استریم شود) و حافظه‌ی اضافی‌ای که گرفتیم هم از 
$O(k)$
است.

}




%%%%%%%%%%%%%%%%%%%%%%%%%%%%%%%%%%%%%%%%%%%%%%%%%%%%%
% برخی دستورات مفید برای تمرین‌ها. %
%%%%%%%%%%%%%%%%%%%%%%%%%%%%%%%%%%%%%%%%%%%%%%%%%%%%%

% فرمول‌های نمونه:
%$f(x) = \Theta(x)$
%$T(x, y) \leq \log(x) + 2^y + \binom{2n}{n}$
% {\tt code\_function}


% لیست موارد بدون شماره‌گذاری:
%\شروع{فقرات}
%    \فقره مورد اول 
%        \شروع{فقرات}
%            \فقره مورد اول داخلی 
%                \شروع{فقرات}
%                    \فقره اولین مورد خیلی داخلی 
%                    \فقره دومین مورد خیلی داخلی 
%                \پایان{فقرات}
%            \فقره مورد دوم داخلی 
%        \پایان{فقرات}
%    \فقره مورد دوم 
%\پایان{فقرات}

% لیست موارد شماره‌گذاری شده:
%\شروع{شمارش}
%    \فقره یک مورد 
%    \فقره یک مورد دیگر 
%    \فقره باز هم یک مورد دیگر
%\پایان{شمارش}

% برای نوشتن یک فرمول چند خطی که یک شماره بخوره. برای حالت‌های دیگر به سایت زیر مراجعه کنید:
% https://www.overleaf.com/learn/latex/Aligning_equations_with_amsmath
%\begin{align} 
%    \begin{split}
%        (x+y)^3 &= (x+y)^2(x+y) \\
%                &= (x^2+2xy+y^2)(x+y) \\
%                &= (x^3+2x^2y+xy^2) + (x^2y+2xy^2+y^3) \\
%                &= x^3+3x^2y+3xy^2+y^3
%    \end{split}                                 
%\end{align}

% نوشتن ماتریس و آرایه دوبعدی
%\begin{align}
%    A = 
%    \begin{bmatrix}
%        A_{11} & A_{21} \\
%        A_{21} & A_{22}
%    \end{bmatrix}
%\end{align}

% اضافه کردن تصویر با شرح. برای اضافه کردن تصویر بدون شرح، فقط دستور درج‌تصویر را استفاده کنید.
%\شروع{شکل}[h!]
%\تنظیم‌ازوسط
%  \درج‌تصویر[پهنا=0.5\پهنای‌متن]{img.jpg}
%\شرح{شرح تصویر}
%\برچسب{تصویر:۱}
%\پایان{شکل}
%\پایان{زیرمسائل}

%کد زیر را ببینید.
%%\lstinputlisting{source_filename.py}
%\begin{latin}
%\begin{lstlisting}
%# Import the modules
%import sys
%import random
%
%ans = True
%
%while ans:
%question = raw_input("Ask the magic 8 ball a question: (press enter to quit) ")
%
%answers = random.randint(1,8)
%
%if question == "":
%sys.exit()
%
%elif answers == 1:
%print "It is certain"
%
%elif answers == 2:
%print "Outlook good"
%
%elif answers == 3:
%print "You may rely on it"
%
%elif answers == 4:
%print "Ask again later"
%
%elif answers == 5:
%print "Concentrate and ask again"
%
%elif answers == 6:
%print "Reply hazy, try again"
%
%elif answers == 7:
%print "My reply is no"
%
%elif answers == 8:
%print "My sources say no"
%\end{lstlisting}
%\end{latin}
